\section{From pair resampling to the switch chain}
\label{mix:section}

We now combine the local estimates, identify the kernel of the resulting
operator, and compare it with the chain in Theorem~\ref{thm:main}.

\subsection{The global operator inequality}

We use the triangle regrouping of Caputo's squared-generator method
\cite[Section~1.1, Equations~(1.13)--(1.15)]{Caputo08}.
The exact correction counts how often each pair belongs to a triple.
Here Corollary~\ref{dis:nonnegative} supplies the disjoint-term
positivity that commutation supplies in that framework.

\begin{proposition}\label{mix:global}
The pair-resampling operator satisfies
\begin{equation}\label{mix:H-square}
H^2\succeq H.
\end{equation}
\end{proposition}

\begin{proof}
Let $D$ be the ordered disjoint-pair sum from
Equation~\eqref{dis:D}. Every ordered pair of distinct intersecting vertex
pairs belongs to the unique triple formed by their union. Each single
vertex pair belongs to $n-2$ triples, and $h_a^2=h_a$. Expanding the
squares therefore gives the exact identity
\begin{equation}\label{mix:expansion}
H^2=\sum_{T\in\binom V3} H_T^2-(n-3)H+D.
\end{equation}
Proposition~\ref{tri:gap} and Corollary~\ref{dis:nonnegative} imply
\[
H^2\succeq\sum_{T\in\binom V3}H_T-(n-3)H
=(n-2)H-(n-3)H=H.
\]
\end{proof}

\subsection{Connectivity and the kernel}

We give a direct connectivity argument to ensure that the operator bound
controls every centered function. It uses Havel's classical
degree-realization switching argument~\cite[Theorem~3]{Havel55},
with deterministic tie-breaking to select the same realization from
every starting graph.

\begin{lemma}\label{mix:connectivity}
Any two graphs in $\Omega_d$ can be joined by a sequence of valid switches.
\end{lemma}

\begin{proof}
We transform every realization into the same recursively prescribed
graph. Suppose a set of vertices has already been frozen, and let $R$ be
the remaining vertex set. The edges involving frozen vertices will no
longer change. Write $d_i^R$ for the degrees in the current induced graph
on $R$.

Choose the first vertex $v\in R$ in label order. Prescribe its desired
neighbors to be the $d_v^R$ vertices of $R\setminus\{v\}$ of highest
residual degree, breaking ties by label. If the current neighbors differ
from this prescribed set, there are a missing desired neighbor $j$ and
an undesired neighbor $i$. Their residual degrees satisfy
$d_j^R\ge d_i^R$. With neighborhoods taken in the induced graph on $R$,
put
\[
A=N_R(j)\setminus\{i\},\qquad
B=N_R(i)\setminus\{j\}.
\]
Removing a possible edge $ij$ subtracts the same amount from both
neighborhood sizes, so $|A|\ge |B|$. Moreover $v\in B\setminus A$.
Consequently $A\setminus B$ is nonempty. A vertex $k$ in this set is
distinct from $v,i,j$, and the membership conditions say precisely that
$vi,jk$ are edges while $vj,ik$ are absent. The valid switch
\[
\{vi,jk\}\longmapsto\{vj,ik\}
\]
changes no other edge at $v$, so it retains every previously attained
desired neighbor and increases their number by one. Repetition therefore
realizes the prescribed neighborhood. If fewer than four residual
vertices remain, a discrepancy would still force the four distinct
vertices just constructed; hence no discrepancy can occur.

Freeze $v$ and continue on $R\setminus\{v\}$. Its residual degree vector
is now determined by the previous vector and the prescribed neighborhood,
independently of the starting realization. It is graphical because the
current graph realizes it. Inductively, the prescription fixes the same
edges from every starting graph. All switches occur entirely among
unfrozen vertices, and the procedure terminates after all vertices are
frozen. Thus every realization is connected to the same graph. Reversing
one of these paths proves the claim.
\end{proof}

\begin{corollary}\label{mix:poincare}
The kernel of $H$ consists exactly of the constant functions. For every
real function $f$ on $\Omega_d$,
\begin{equation}\label{mix:H-poincare}
\operatorname{Var}_{\pi_d}(f)\le \langle f,Hf\rangle.
\end{equation}
\end{corollary}

\begin{proof}
Since the $h_a$ are orthogonal projections,
\[
\langle f,Hf\rangle=\sum_{a\in\binom V2}\|h_af\|^2.
\]
Its vanishing forces $f$ to be constant on every pair fiber. Every valid
switch exchanges two singleton assignments in a pair fiber: take either
opposite vertex pair of the toggled four-cycle. Thus $f$ is constant along
switches, and Lemma~\ref{mix:connectivity} makes it constant on $\Omega_d$.
The converse is immediate.

The operator $H$ is positive semidefinite. By
Proposition~\ref{mix:global}, each eigenvalue $\lambda$ satisfies
$\lambda^2\ge\lambda$, and hence is either zero or at least one.
Its kernel has just been identified, so its restriction to the centered
subspace is bounded below by the identity. Subtracting the mean of $f$
gives Equation~\eqref{mix:H-poincare}.
\end{proof}

\subsection{Comparison with single switches}

We compare a whole-fiber resampling with exchanges of two singleton
assignments inside that fiber. This is the fiberwise heat-bath
comparison principle used by Carstens and Kleer for bipartite switch
and Curveball chains~\cite[Theorems~2 and~7 and Appendix~A]{CK18}.
The argument below proves the required undirected comparison with
its own constant and counts the multiplicity of each switch.

Let $L$ be the unnormalized Laplacian of the switch graph:
\begin{equation}\label{mix:L}
(Lf)(G)=\sum_{G'\text{ a switch neighbor of }G}
\bigl(f(G)-f(G')\bigr).
\end{equation}
Each distinct neighbor appears once in this sum.

For a vertex pair $a$, define $J_a$ on each of its fibers by
\[
J_a=\sum_{1\le k<\ell\le N}(I-\tau_{k\ell}),
\]
where the $N$ singleton vertices have been indexed and $\tau_{k\ell}$
exchanges their assignments. Combining these operators over all fibers
defines $J_a$ on the full graph space.

\begin{lemma}\label{mix:comparison}
The operators above satisfy
\begin{equation}\label{mix:comparison-bounds}
h_a\preceq n^2J_a,\qquad
\sum_{a\in\binom V2}J_a=2L,\qquad
H\preceq2n^2L.
\end{equation}
\end{lemma}

\begin{proof}
First work in one nonconstant uniform $q$-subset fiber with $N$ singleton
labels. Let $S_0$ be a uniform $q$-subset and let $\sigma$ be an independent
uniform permutation of the $N$ labels. Then $\sigma S_0$ is uniform and
independent of $S_0$, so
\[
\operatorname{Var}(f)=\frac12\mathbb E
\bigl(f(S_0)-f(\sigma S_0)\bigr)^2.
\]
Fix, for each permutation, a decomposition into at most $N$
transpositions, obtained for example by decomposing its cycles. Pad with
identities to length $N$. Conditional on $\sigma$, every intermediate
subset is uniform. Cauchy--Schwarz along the resulting path gives
\begin{align*}
\operatorname{Var}(f)
&\le \frac{N^2}{2}\sum_{k<\ell}
\mathbb E_{\mathrm{slice}}
\bigl(f-\tau_{k\ell}f\bigr)^2\\
&=N^2\left\langle f,
\sum_{k<\ell}(I-\tau_{k\ell})f\right\rangle.
\end{align*}
Indeed, Cauchy--Schwarz costs a factor $N$, and each of the $N$ steps has
conditional squared difference bounded by the sum of all transposition
squared differences displayed above. Identity steps contribute zero.
Constant fibers require no estimate. Since $N\le n$, averaging over
fibers proves $h_a\preceq n^2J_a$.

Each switch has two pair-trade realizations, as observed in
Lemma~2 of~\cite{CBS18}. To check this count in the present notation,
consider a particular transition
\[
\{uv,xy\}\longmapsto\{ux,vy\}.
\]
The four distinct vertices form the toggled cycle $u,v,y,x,u$. Taking
$a=\{u,y\}$ exchanges the assignments of singleton vertices $v,x$;
taking $a=\{v,x\}$ exchanges those of singleton vertices $u,y$. These
are two occurrences of this transition. They are also the only ones:
a singleton transposition toggles precisely the four edges between its
indexing pair and the two exchanged singleton vertices. Thus its
indexing pair must be one of the two parts of the unique bipartition of
the toggled four-cycle. Each choice determines its transposition uniquely.
It follows pointwise that $\sum_aJ_a=2L$. Summing the first comparison
now proves the last one.
\end{proof}

Combining Corollary~\ref{mix:poincare} and Lemma~\ref{mix:comparison}, we
obtain
\begin{equation}\label{mix:L-poincare}
\operatorname{Var}_{\pi_d}(f)\le 2n^2\langle f,Lf\rangle.
\end{equation}
In particular, if $|\Omega_d|>1$, the spectral gap of $L$ is at least
$1/(2n^2)$.

\subsection{The prescribed chain and its mixing time}

\begin{proof}[Proof of Theorem~\ref{thm:main}]
If $|\Omega_d|=1$, the mixing time is zero. Suppose henceforth that
$|\Omega_d|>1$. For any distinct switch neighbors $G,G'$, their transition
uniquely specifies its four-element vertex set and the ordered pair of
removed and added matchings. Under the proposal rule of
Section~\ref{sec:introduction},
\[
P_d(G,G')=\frac12\,\frac1{\binom n4}\,\frac16
=\frac1{12\binom n4}.
\]
The reverse transition has the same probability, so $\pi_d$ is
reversible, and the diagonal probabilities give
\begin{equation}\label{mix:normalization}
I-P_d=\frac{L}{12\binom n4}.
\end{equation}
By Equation~\eqref{mix:L-poincare}, the spectral gap of $P_d$ is at least
\begin{equation}\label{mix:gamma}
\gamma_0=\frac1{24n^2\binom n4}.
\end{equation}

The holding probability also ensures nonnegative eigenvalues: write
$P_d=(I+Q)/2$, where $Q$ is the reversible Markov kernel consisting of
the proposal step, including rejected proposals. Jensen's inequality and
stationarity show that $Q$ is a contraction on $L^2(\pi_d)$, so its
eigenvalues lie in $[-1,1]$. Thus every eigenvalue of $P_d$ on the centered
subspace lies in $[0,1-\gamma_0]$, and for a centered function $f$,
\[
\|P_d^t f\|\le(1-\gamma_0)^t\|f\|
\le e^{-\gamma_0t}\|f\|.
\]
The density of a point mass at $G$, minus one, has squared norm
$|\Omega_d|-1$. Reversibility evolves this density by $P_d$. Therefore
Cauchy--Schwarz yields, uniformly over starting states,
\begin{equation}\label{mix:TV}
\|P_d^t(G,\cdot)-\pi_d\|_{\mathrm{TV}}
\le\frac12\sqrt{|\Omega_d|}\,e^{-\gamma_0t}.
\end{equation}
In particular, total variation is at most $1/4$ when
\[
t=\left\lceil
24n^2\binom n4\left(\log2+\frac12\log|\Omega_d|\right)
\right\rceil.
\]
Here and throughout, logarithms are natural. Since there are at most
$2^{\binom n2}$ labeled simple graphs,
\[
\log2+\frac12\log|\Omega_d|
\le \left(1+\frac{n(n-1)}4\right)\log2\le n^2
\qquad(n\ge4).
\]
Also $24n^2\binom n4\le n^6$. The displayed time is consequently at most
$n^8+1\le2n^8$, as claimed.
\end{proof}
